\section*{参考文献}
\begingroup
\scriptsize
\setlength{\parindent}{0pt}
\setlength{\parskip}{1.4pt}
\noindent 下列 12 个来源被正文实际引用；归档共 14 个。完整标题见 \texttt{sources.csv}。\par
[1] ServerBench: How can your coding agents de\ldots{}\hfill{\texttt{2611.06841}}\par
[2] VibeServe: Can AI Agents Build Bespoke LLM\ldots{}\hfill{\texttt{2605.06068}}\par
[3] SEIS: Self-Evolving Inference Systems\hfill{\texttt{2610.04646}}\par
[4] InferenceBench: A Benchmark for Open-Ended\ldots{}\hfill{\texttt{2607.20468}}\par
[5] FlashRT: Agent Harness for Guiding Agents\ldots{}\hfill{\texttt{2607.18171}}\par
[6] Simthesizer: An Agent-Driven Simulation Fr\ldots{}\hfill{\texttt{2608.24650}}\par
[7] ISO-Bench: Can Coding Agents Optimize Real\ldots{}\hfill{\texttt{2602.19594}}\par
[8] SWE-Serve: Benchmarking Agentic Engineerin\ldots{}\hfill{\texttt{2609.26777}}\par
[9] InfraDev: Can Coding Agents Develop AI Inf\ldots{}\hfill{\texttt{2610.36589}}\par
[10] AutoTuneBench: Trustworthy Measurement for\ldots{}\hfill{\texttt{2609.18123}}\par
[11] KernelBench: Can LLMs Write Efficient GPU\ldots{}\hfill{\texttt{2502.10517}}\par
[12] KernelBench-Verified: Do LLM-Generated Ker\ldots{}\hfill{\texttt{2607.16241}}\par
\endgroup
